\documentclass{article}
\usepackage{nips15submit_e,times}
\usepackage{hyperref}
\usepackage{url}


\title{Learning to Teach}


\author{
Gordon Chen\\
School of Computing\\
University of Connecticut\\
\texttt{gbc24001@uconn.edu}\\
}

\newcommand{\fix}{\marginpar{FIX}}
\newcommand{\new}{\marginpar{NEW}}

\nipsfinalcopy

\begin{document}


\maketitle

\begin{abstract}
\end{abstract}

\section{Problem}
LLMs today are not native teachers. They give out answers too quickly, give explanations that are
miscalibrated for the user's knowledge level, and in general serve the role of an encyclopedia more
than that of a teacher who knows how to help a student along.

While it's not fair to blame LLMs for lazy students just getting the answer from LLMs, it is a problem
that we as model developers can address when students that want to learn are deprived of the ability
to struggle through a problem when an LLM gives them the solution immediately instead of nudging them
in the right direction.

\section{Goal}
The goal of this project is to train an LLM to be the best possible teacher, to natively optimize
for teaching ability so that students can benefit the most from trying to learn from LLMs.

Essentially I want use reinforcement learning to fine-tune a teacher LLM to provide the most useful
feedback for a student LLM. The student LLM would be frozen, but would be able to ``learn" in-context from
the teacher's feedback. For every problem a student LLM gets wrong, the teacher LLM would be given
the question and the student LLM's response, and ouput feedback for the student. That feedback is then placed
in the student's context, and the student is evaulated on similar questions generated by frozen,
question-writing LLM. As a baseline, the student would also be evaulated on those questions without
the teacher's feedback in its context, and the difference in the student's accuracy would serve as a learning
signal for the teacher to know if its feedback is helpful or not.

By evaluating the student on similar problems the teacher has never seen, we prevent the teacher from just
giving the student the answer and calling it a day because the teacher's answer must help the student
generalize to the similar questions.

An additional feature I could explore if I have time would be multi-turn interactions between the teacher
and the student. Instead of just giving feedback a single time, the student would generate a solution,
the teacher would give feedback, the student would do a similar problem with the feedback, and repeat.

\section{Importance}
This is an important project because as LLMs are getting smarter, their potential to become
really cheap, widely avaiable, infinitely-patient, competent personal tutors is increasing.
However, currently many students are using LLMs in a very destructive way, using LLMs to offload
the need to think and learn by just asking LLMs to do their homework. Not just that, but for all of the
work going into making LLMs smarter, in my opinion, not enough focus has been put into making the information
transfer between LLMs (which are getting smarter very quickly) and humans better. I hope that this project
shows that reinforcement learning can teach LLMs to be really good teachers, and that there is a lot of
low hanging fruit in terms of making them effective teachers.

\section{Expected Results}

\section{Data}
I plan on synthetically generating a dataset with multiple concepts and multiple questions per concept.
While this paradigm of learning to teach could potentially work on any subject where it is possible to
objectively know if a student is doing better (multiple choice is posssible for English and general knowledge),
I plan on training a teacher on math. The dataset would then have multiple relatively simple math techniques,
such as solving quadratic equations, math word problems, and solving systems of equations, that I can hopefully
easily synthetically generate, create similar but diverse enough variations of, and be relatively easy to verify.

Since the goal of this project is to see if a teacher can be trained to improve student performance,
these synthetically generated problems need to be hard enough that student performance is initally bad,
but easy enough that with the proper teacher guidance, it is possible for a small ,since I am compute
constrained, student model to find the solution.

\section{Algorithms}
I plan to use GRPO and LoRA to fine-tune the teacher LLM. By using LoRA and fine-tuning a model instead
of training a model from scratch I hope to make this project feasible to develop on my RTX 3060, and if not,
at least not cost too many cloud compute credits.

While I am only fine-tuning the teacher LLM, I also need to be able to sample from a (potentially smaller)
student LLM, as well as a question writing model, although perhaps some problems can be programatically
generated with substitution instead of with another LLM.

\section{Schedule and Milestones}

\subsubsection*{Acknowledgments}
\subsubsection*{References}
\end{document}
